\documentclass[12pt,a4paper]{article}

\usepackage[utf8]{inputenc}
\usepackage[T1]{fontenc}
\usepackage[romanian]{babel}
\usepackage[
    a4paper,
    top=1.5cm,
    bottom=1.9cm,
    left=1.2cm,         % Marginea din stânga
    right=1.2cm         % Marginea din dreapta
]{geometry}
\usepackage{listings}
\usepackage{enumitem}
\usepackage{xcolor}
\usepackage{amssymb}
\usepackage{multicol}
\usepackage{titlesec} % Pachet esențial pentru spațierea titlurilor
\usepackage{parskip} % Pachet pentru controlul spațiului între paragrafe

% --- HYPERREF SETUP ---
% Pachetul pentru link-uri, acum cu opțiuni
\usepackage[colorlinks=true, linkcolor=blue, urlcolor=blue]{hyperref}

% --- SETARI SPATIERE ---
\definecolor{codegreen}{rgb}{0,0.6,0}
\definecolor{codegray}{rgb}{0.5,0.5,0.5}
\definecolor{codepurple}{rgb}{0.58,0,0.82}
\definecolor{backcolour}{rgb}{0.98,0.98,0.98}
\setlist[enumerate]{nosep} % Setează listele enumerate să nu aibă spațiu vertical

% Setează spațierea pentru titluri
\titlespacing*{\section}{0pt}{2ex}{1ex}
\titlespacing*{\subsection}{0pt}{1.5ex}{0.5ex} % Foarte strâns pentru subsecțiuni (întrebări)
\setlength{\parskip}{0pt} % Elimină spațiul dintre paragrafe

% --- LISTINGS SETUP ---
\lstdefinestyle{mystyle}{
    backgroundcolor=\color{backcolour},   
    commentstyle=\color{codegreen},
    keywordstyle=\color{blue},
    numberstyle=\tiny\color{codegray},
    stringstyle=\color{codepurple},
    basicstyle=\ttfamily\footnotesize,
    breakatwhitespace=false,         
    breaklines=true,                 
    captionpos=b,                    
    keepspaces=true,                 
    showspaces=false,                
    showstringspaces=false,
    showtabs=false,                  
    tabsize=2,
    aboveskip=0ex, % Spatiu redus inainte de listing
    belowskip=0ex  % Spatiu redus dupa listing
}
\lstset{style=mystyle}

\lstdefinelanguage{x86asm}{
  keywords={mov, movl, movb, add, addl, sub, subl, mul, imul, div, divl, idivl, cdq, int, jmp, je, jne, jg, jl, jge, jle, cmp, cmpl, lea, leal, loop, inc, dec, shr, shl, sal, push, pop, call, ret, .global, .text, .data, .asciz, .ascii, .long, .word, .byte, .space},
  morecomment=[l]{\#}, 
  sensitive=false
}

% --- TITLU ---
\title{Compilație Probleme Laborator ASC \\ (x86 AT\&T Syntax)}
\author{}
\date{}
\setcounter{tocdepth}{2}

\begin{document}
\maketitle
\tableofcontents
\newpage

% \setlength{\columnsep}{0.5cm}
% \begin{multicols}{2}

\section{Test Laborator 2.1 / 2.2 / 2.3} \label{sec:lab2} \hfill \hyperref[sec:rezolvari]{\small (vezi răspunsuri)}
% --- Întrebarea 1 (din 2.1/2.2) ---
\subsection*{Întrebarea 1}
Ce valoare va reține EAX după executarea următoarei secvențe de instrucțiuni?

\begin{lstlisting}[language=x86asm]
movl $0, %eax
movb $4, %ah
movb $2, %al
\end{lstlisting}

\textbf{Opțiuni:}
\begin{multicols}{2}
\begin{enumerate}[label=\alph*)]
    \item 513
    \item 1026
    \item 0
    \item 258
\end{enumerate}
\end{multicols}

% --- Întrebarea 2 (din 2.1/2.2) ---
\subsection*{Întrebarea 2}
Se consideră declarate \texttt{x: .word 1} și \texttt{y: .word 2} (sau \texttt{x: .word 1, y: .word 0, z: .word 2}). Ce valoare va avea eax după executarea instrucțiunii \texttt{mov x, \%eax}?

\textbf{Opțiuni:}
\begin{multicols}{2}
\begin{enumerate}[label=\alph*)]
    \item 1 
    \item 2
    \item 0x00020001
\end{enumerate}
\end{multicols}

% --- Întrebarea 3 (din 2.1/2.2) ---
\subsection*{Întrebarea 3}
Fie următoarea declarare în secțiunea .data:
\begin{lstlisting}[language=x86asm]
str1: .ascii "abc"
str2: .ascii "123" # (sau .ascii "1")
\end{lstlisting}


Ce se va afișa în urma apelului WRITE următor?
\begin{lstlisting}[language=x86asm]
movl $4, %eax
movl $1, %ebx
movl $str1, %ecx
movl $5, %edx  # (sau $4, %edx)
int $0x80
\end{lstlisting}

\begin{multicols}{2}
\textbf{Opțiuni (pentru testul cu "123" și edx=5):}
\begin{enumerate}[label=\alph*)]
    \item abc12 
    \item abc + \textit{o valoare reziduală}
    \item nimic
    \item abc
\end{enumerate}

\textbf{Opțiuni (pentru testul cu "1" și edx=4):}
\begin{enumerate}[label=\alph*)]
    \item nimic
    \item abc1 
    \item abc
    \item abc + \textit{o valoare reziduală}
\end{enumerate}
\end{multicols}

% --- Întrebarea 4 (din 2.1/2.2) ---
\subsection*{Întrebarea 4}
Fie programul următor. Ce valoare va fi afișată în \textbf{\%EAX} în urma rulării comenzilor \texttt{b main}, \texttt{run}, \texttt{stepi}, \texttt{stepi}, \texttt{i r}?

\begin{lstlisting}[language=x86asm]
.data
    x: .long 0x04030201
    y: .long 0x08070605
.text
.global main
main:
    mov x, %eax
    mov y, %al  # (sau %ah)
    mov $1, %eax
    mov $0, %ebx
    int $0x80
\end{lstlisting}

\begin{multicols}{2}
\textbf{Opțiuni (pentru \texttt{mov y, \%al} sau \texttt{mov y, \%ah}):}
\begin{enumerate}[label=\alph*)]
    \item 0x08070605
    \item 0x04030205 
    \item 0x04030501
    \item 1
    \item 0x04030201
\end{enumerate}
\end{multicols}

% --- Întrebarea 5 (din 2.3) ---
\subsection*{Întrebarea 5}
Ce valoare va reține registrul CH după executarea următoarei instrucțiuni?
\begin{lstlisting}[language=x86asm]
movl $553, %ecx
\end{lstlisting}

\textbf{Opțiuni:}
\begin{multicols}{2}
\begin{enumerate}[label=\alph*)]
    \item 0
    \item 2 
    \item 10
    \item 512
\end{enumerate}
\end{multicols}

% --- Întrebarea 6 (din 2.3) ---
\subsection*{Întrebarea 6}
Fie următoarea declarare în secțiunea .data:
\begin{lstlisting}[language=x86asm]
str: .ascii "1234"
x: .byte 97
\end{lstlisting}
Ce se va afișa în urma apelului WRITE următor?
\begin{lstlisting}[language=x86asm]
movl $4, %eax
movl $1, %ebx
movl $str, %ecx
movl $5, %edx
int $0x80
\end{lstlisting}

\textbf{Opțiuni:}
\begin{multicols}{2}
\begin{enumerate}[label=\alph*)]
    \item 1234
    \item 1234a
    \item 12349
    \item 123497
\end{enumerate}
\end{multicols}

% --- Întrebarea 7 (din 2.3) ---
\subsection*{Întrebarea 7}
Fie următorul program. Precizați secvența corectă de instrucțiuni în debugger, în urma căreia vom obține valoarea 8.
\begin{lstlisting}[language=x86asm]
.data
.text
.global main
main:
    movl $8, %eax
    movl $2048, %ecx
    et_exit:
    movl $1, %eax
    movl $0, %ebx
    int $0x80
\end{lstlisting}

\textbf{Opțiuni:}
\begin{enumerate}[label=\alph*)]
    \item \texttt{b main; run; i r eax}
    \item \texttt{b main; run; stepi; stepi; i r cl}
    \item \texttt{b main; run; stepi; stepi; i r ch} 
    \item \texttt{b main; run; stepi; i r ah}
\end{enumerate}
\newpage
\section{Test Laborator 3.1 / 3.2} \label{sec:lab3} \hfill \hyperref[sec:rezolvari]{\small (vezi răspunsuri)}

% --- Întrebarea 1 (din 3.1) ---
\subsection*{Întrebarea 1}
Fie codul de mai jos. Care este valoarea lui s c\^and execu\c{t}ia ajunge la \texttt{et\_exit}?
\begin{multicols}{2}
\begin{lstlisting}[language=x86asm]
.data
    s: .long 0
.text
.global main
main:
    mov $1, %edx
    mov $0, %eax
    movl $0xffffffff, %ebx
    divl %ebx
    mov %eax, %ecx

et_loop:
    add %ecx, s
    loop et_loop

et_exit:
    mov $1, %eax
    mov $0, %ebx
    int $0x80
\end{lstlisting}
\end{multicols}

\textbf{Opțiuni:}
\begin{multicols}{2}
\begin{enumerate}[label=\alph*)]
    \item 1
    \item 0
    \item 0xffffffff
    \item 15
\end{enumerate}
\end{multicols}

% --- Întrebarea 2 (din 3.1) ---
\subsection*{Întrebarea 2}
Se stochează în registrul eax valoarea \texttt{0x40000000}, în ebx \texttt{0x8}, în ecx \texttt{0x1} și în edx \texttt{0x8}. Ce valori vor avea registrii eax si edx dupa executarea instructiunii \texttt{mul \%edx}?

\textbf{Opțiuni:}
\begin{multicols}{2}     
\begin{enumerate}[label=\alph*)]
    \item eax=32, edx=0
    \item eax=0, edx=2
    \item eax=2, edx=0
    \item eax=32, edx=2
\end{enumerate}
\end{multicols}

% --- Întrebarea 3 (din 3.1) ---
\subsection*{Întrebarea 3}
Se stochează în \%edx valoarea 0, în eax 47 și în ebx 15. Ce valori vor avea registrii eax si edx dupa executarea instructiunii \texttt{div \%ebx}?

\textbf{Opțiuni:}
\begin{multicols}{2}     
\begin{enumerate}[label=\alph*)]
    \item eax=2, edx=3
    \item eax=3, edx=0
    \item eax=3, edx=2
    \item eax=2, edx=0
\end{enumerate}
\end{multicols}

% --- Întrebarea 4 (din 3.1) ---
\subsection*{Întrebarea 4}
Fie codul de mai jos. Care este valoarea depozitată la final în ecx (în dreptul etichetei \texttt{exit})?

\begin{multicols}{2}
\begin{lstlisting}[language=x86asm]
.data
    x: .long 0x80000000
    y: .long 0x70000000
.text
.global main
main:
    mov x, %eax
    cmp y, %eax
    jge label
    mov $5, %ecx
    jmp exit
label:
    mov $6, %ecx
exit:
    mov $1, %eax
    mov $0, %ebx
    int $0x80
\end{lstlisting}
\end{multicols}

\textbf{Opțiuni:}
\begin{multicols}{2}
\begin{enumerate}[label=\alph*)]
    \item 5
    \item 6
\end{enumerate}
\end{multicols}
\newpage

% --- Întrebarea 5 (din 3.1) ---
\subsection*{Întrebarea 5}
Fie următorul program. Ce valoare vom obține dacă vom rula cu debuggerul \texttt{b et\_exit}, \texttt{run}, \texttt{ir ebx}?
\begin{multicols}{2}
\begin{lstlisting}[language=x86asm]
.data
.text
.global main
main:
    mov $3, %eax
    shl $2, %eax
    mov $2, %ebx
    mul %ebx
    mov $0, %edx
    mov $8, %ebx
    div %ebx
    sub %eax, %ebx
et_exit:
    mov $1, %eax
    mov $0, %ebx
    int $0x80
\end{lstlisting}
\end{multicols}

\textbf{Opțiuni:}
\begin{multicols}{2}
\begin{enumerate}[label=\alph*)]
    \item 8
    \item 3
    \item 5
    \item 2
\end{enumerate}
\end{multicols}

% --- Întrebarea 6 (din 3.1) ---
\subsection*{Întrebarea 6}
Care este ordinea de trecere prin etichete?

\begin{multicols}{2}
\begin{lstlisting}[language=x86asm]
.data
.text
.global main
main:
    mov $1, %eax
    mov $2, %ebx
    mov $3, %ecx
    mov $4, %edx
    cmp %ebx, %eax
    je etx
ety:
    cmp %ecx, %edx
    jg etz
    jmp ett
etx:
    jmp ety
etz:
    mov %ebx, %edx
    jmp ety
ett:
    mov $1, %eax
    mov $0, %ebx
    int $0x80
\end{lstlisting}
\end{multicols}

\textbf{Opțiuni:}
\begin{multicols}{2}
\begin{enumerate}[label=\alph*)]
    \item ety, etz, ety, ett
    \item etx, ety, ett
    \item ety, etx, etz, ett
    \item etx, ety, etz, ety, ett
\end{enumerate}
\end{multicols}

% --- Întrebarea 7 (din 3.2) ---
\subsection*{Întrebarea 7}
Fie codul de mai jos. Care sunt valorile lui eax și edx când execuția ajunge la \texttt{label}?
\begin{multicols}{2}
\begin{lstlisting}[language=x86asm]
.data
    x: .long 17
    y: .long 6
.test
.global main
main:
    mov $1, %edx
    mov x, %eax
    jmp et
    mov $0, %edx
et:
    divl y
label:
    mov $1, %eax
    mov $0, %ebx
    int $0x80
\end{lstlisting}     
\end{multicols}

\textbf{Opțiuni:}
\begin{multicols}{2}
\begin{enumerate}[label=\alph*)]
    \item eax=0x2, edx=0x5
    \item eax=0x2aaaaaad, edx=0x3
    \item eax=0x5, edx=0x2
    \item eax=0x3, edx=0x2aaaaaad
\end{enumerate}
\end{multicols}
\newpage

% --- Întrebarea 8 (din 3.2) ---
\subsection*{Întrebarea 8}
Se stochează în EAX valoarea \texttt{0x80000000}, în EBX \texttt{0x8}, în ECX \texttt{0x1} și în EDX \texttt{0x4}. Ce valori vor avea registrii EAX, respectiv EDX după executarea instrucțiunii \texttt{mul \%ebx}?

\textbf{Opțiuni:}
\begin{multicols}{2}
\begin{enumerate}[label=\alph*)]
    \item eax=0, edx=4
    \item eax=32, edx=0
    \item eax=4, edx=0
    \item eax=32, edx=4
\end{enumerate}
\end{multicols}

% --- Întrebarea 9 (din 3.2) ---
\subsection*{Întrebarea 9}
Se stochează în \%edx valoarea 0, în eax 37 și în ebx 15. Ce valori vor avea registrii eax si edx dupa executarea instructiunii \texttt{div \%ebx}?

\textbf{Opțiuni:}
\begin{multicols}{2}
\begin{enumerate}[label=\alph*)]
    \item eax=7, edx=0
    \item eax=2, edx=7
    \item eax=2, edx=0
    \item eax=7, edx=2
\end{enumerate}
\end{multicols}

% --- Întrebarea 10 (din 3.2) ---
\subsection*{Întrebarea 10}
Fie codul de mai jos. Care este valoarea depozitată în z când ajungem la eticheta \texttt{final}?

\begin{multicols}{2}
\begin{lstlisting}[language=x86asm]
.data
    x: .long 17
    y: .long 6
    x1: .long 5
    y1: .long 9
    z: .space 4
.text
.global main
main:
    mov x, %eax
    mov y, %ebx
    cmp %eax, %ebx
    jge et1
    mov x1, %eax
    cmp %eax, %ebx
    jle et
    mov x, %ebx
et:
    mov x1, %eax
    mov y1, %ebx
    cmp %eax, %ebx
    jge et2
    add %eax, %ebx
    jmp final
et1:
    add %eax, %ebx
    jmp final
et2:
    sub %eax, %ebx
final:
    mov %ebx, z
    mov $1, %eax
    mov $0, %ebx
    int $0x80
\end{lstlisting}
\end{multicols}

\textbf{Opțiuni:}
\begin{multicols}{2}
\begin{enumerate}[label=\alph*)]
    \item 12
    \item 4 
    \item 1
    \item 23
\end{enumerate}
\end{multicols}
\newpage

% --- Întrebarea 11 (din 3.2) ---
\subsection*{Întrebarea 11}
Care este ordinea de trecere prin etichete?

\begin{multicols}{2}
\begin{lstlisting}[language=x86asm]
.data
.text
.global main
main:
    jmp etb
eto:
    jmp etd
eth:
    jmp eto
etb:
    jmp eth
etd:
    mov $1, %eax
    mov $0, %ebx
    int $0x80
\end{lstlisting}
\end{multicols}

\textbf{Opțiuni:}
\begin{multicols}{2}
\begin{enumerate}[label=\alph*)]
    \item etb, eto, eth, etd
    \item etb, eth, eto, etd 
    \item eto, eth, etb,etd
    \item eth, etb, etd, eto
\end{enumerate}
\end{multicols}

% --- Întrebarea 12 (din 3.2) ---
\subsection*{Întrebarea 12}
De câte ori se va executa instrucțiunea \texttt{loop}?
\begin{multicols}{2}
\begin{lstlisting}[language=x86asm]
.data
    x: .long 5
    y: .long 5
    s: .long 0
.text
.global main
main:
    mov x, %ecx
    sub y, %ecx
et:
    add %ecx, s
    loop et
exit:
    mov $1, %eax
    mov $0, %ebx
    int $0x80
\end{lstlisting}
\end{multicols}

\textbf{Opțiuni:}
\begin{multicols}{2}
\begin{enumerate}[label=\alph*)]
    \item 0x5
    \item 0x1
    \item este un ciclu infinit
    \item 0x0 
    \item 0xffffffff
\end{enumerate}
\end{multicols}

% --- Întrebarea 13 (din 3.2) ---
\subsection*{Întrebarea 13}
Fie următorul program. Ce valoare vom obține dacă vom rula cu debuggerul \texttt{b et\_exit}, \texttt{run}, \texttt{ir edx}?
\begin{lstlisting}[language=x86asm]
.data
.text
.global main
main:
    mov $2, %eax
    mov $3, %ebx
    add %eax, %ebx
    mul %ebx
    mov $0, %edx
    divl $3
    add %eax, %edx
et_exit:
    mov $1, %eax
    mov $0, %ebx
    int $0x80
\end{lstlisting}

\textbf{Opțiuni:}
\begin{multicols}{2}
\begin{enumerate}[label=\alph*)]
    \item 0x4 
    \item 0x3
    \item 0x1
    \item 0x0
\end{enumerate}
\end{multicols}

\section{Test Laborator 4.1 / 4.2 / 4.3} \label{sec:lab4} \hfill \hyperref[sec:rezolvari]{\small (vezi răspunsuri)}
% --- Întrebarea 1 (din 4.1 / 4.2 / 4.3) ---
\subsection*{Întrebarea 1}
Ce valori vor fi depozitate în v când execuția va ajunge în dreptul etichetei \texttt{et\_exit}?

\begin{multicols}{2}
\begin{lstlisting}[language=x86asm]
.data
    v: .space 20  # (sau .space 24)
    n: .long 5
.text
.global main
main:
    lea v, %edi  # (sau movl $v, %edi)
    mov $11, %edx
    mov $0, %ecx

et_loop:
    cmp n, %ecx
    jg et_exit
    mov %edx, (%edi, %ecx, 4) 
    inc %ecx
    inc %edx
    jmp et_loop
et_exit:
    mov $1, %eax
    xor %ebx, %ebx
    int $0x80
\end{lstlisting}
\end{multicols}

\textbf{Opțiuni:}
\begin{multicols}{2}
\begin{enumerate}[label=\alph*)]
    \item 0, 1, 2, 3, 4, 5
    \item 11, 12, 13, 14, 15, 16 
    \item 11, 12, 13, 14, 15
    \item 11, 12, ..., 27
    \item Execuția nu ajunge la et\_exit
\end{enumerate}
\end{multicols}

% --- Întrebarea 2 (din 4.1) ---
\subsection*{Întrebarea 2}
Ce se va afișa pe ecran? (Notă: codul a fost reasamblat din fragmente)

\begin{multicols}{2}
\begin{lstlisting}[language=x86asm]
.data
    x: .long 1, 3, 6, 7, 9
    n1: .long 5
    n2: .long 10
    c: .long 0x64636261
    s: .space 11
.text
.global main
main:
    mov $s, %edi
    mov $x, %esi
    movb c, %al
    mov $0, %ecx
et_loop1:
    cmp n2, %ecx
    je et_exit_loop1
    mov %al, (%edi, %ecx, 1)
    inc %ecx
    jmp et_loop1
et_exit_loop1:
    mov $c, %eax
    movb 1(%eax), %al
    mov $0, %ecx
et_loop2:
    cmp n1, %ecx
    je et_exit
    mov (%esi, %ecx, 4), %ebx
    mov %al, (%edi, %ebx, 1)
    inc %ecx
    jmp et_loop2
et_exit:
    mov $10, %ecx
    movb $0, (%edi, %ecx, 1)
    mov $4, %eax
    mov $1, %ebx
    mov $s, %ecx
    mov $11, %edx
    int $0x80
    mov $1, %eax
    xor %ebx, %ebx
    int $0x80
\end{lstlisting}
\end{multicols}

\textbf{Opțiuni:}
\begin{multicols}{2}
\begin{enumerate}[label=\alph*)]
    \item ababaabbab 
    \item aaaaaaaaaa
    \item bbbbbaaaaa
    \item 61, 61, 61, 61, 61, 61, 61, 61, 61, 61
\end{enumerate}
\end{multicols}
\newpage

% --- Întrebarea 3 (din 4.1) ---
\subsection*{Întrebarea 3}
Ce se va afișa pe ecran?

\begin{multicols}{2}
\begin{lstlisting}[language=x86asm]
.data
    n: .long 3
    s: .asciz "abc"
.text
.global main
main:
    mov $s, %edi
    mov $0, %ecx
et_loop:
    cmp n, %ecx
    je et_exit
    mov (%edi, %ecx, 1), %al
    sub $'a', %al
    add $'A', %al
    mov %al, (%edi, %ecx, 1)
    inc %ecx
    jmp et_loop
et_exit:
    mov $4, %eax
    mov $1, %ebx
    mov $s, %ecx
    mov $4, %edx
    int $0x80
    mov $1, %eax
    xor %ebx, %ebx
    int $0x80
\end{lstlisting}
\end{multicols}

\textbf{Opțiuni:}
\begin{multicols}{2}
\begin{enumerate}[label=\alph*)]
    \item Abc
    \item ABC 
    \item abc
    \item Abc + o valoarea reziduala
\end{enumerate}
\end{multicols}

% --- Întrebarea 4 (din 4.1) ---
\subsection*{Întrebarea 4}
Ce se va afișa pe ecran?
\begin{multicols}{2}
\begin{lstlisting}[language=x86asm]
.data
    n: .long 3
    s: .byte 'a', 'b', 'c'
    t: .byte 'd', 'e', 'f'
    u: .space 4
.text
.global main
main:
    mov $0, %ecx
et_loop:
    cmp n, %ecx
    je et_exit
    mov $0, %edx
    sub %ecx, %edx
    mov t(, %edx, 1), %al
    mov %al, u(, %ecx, 1)
    inc %ecx
    jmp et_loop
et_exit:
    movb $0, u(, %ecx, 1)
    mov $4, %eax
    mov $1, %ebx
    mov $u, %ecx
    mov $4, %edx
    int $0x80
    mov $1, %eax
    xor %ebx, %ebx
    int $0x80
\end{lstlisting}
\end{multicols}

\textbf{Opțiuni:}
\begin{multicols}{2}
\begin{enumerate}[label=\alph*)]
    \item abc
    \item cba
    \item def
    \item dcb 
\end{enumerate}
\end{multicols}
\newpage

% --- Întrebarea 5 (din 4.1) ---
\subsection*{Întrebarea 5}
Fie următorul program. Ce valoare vom obține dacă vom rula \texttt{b et\_exit}, \texttt{run}, \texttt{ir eax}?
\begin{lstlisting}[language=x86asm]
.data
    n: .long 4
    v: .long 0x01020304, 0x05060708, 0x090a0b0c, 0x0d0e0f10
.text
.global main
main:
    mov $v, %esi
    mov $1, %ecx
    mov (%esi, %ecx, 4), %eax
    add $4, %esi
    movb (%esi, %ecx, 4), %al
et_exit:
    mov $1, %eax
    xor %ebx, %ebx
    int $0x80
\end{lstlisting}

\textbf{Opțiuni:}
\begin{multicols}{2}
\begin{enumerate}[label=\alph*)]
    \item 0x0506070c 
    \item 0x090a0b08
    \item 0x05060704
    \item 0x090a0b0c
\end{enumerate}
\end{multicols}

% --- Întrebarea 6 (din 4.1 / 4.3) ---
\subsection*{Întrebarea 6}
Fie următorul program. Ce valoare vom obține dacă vom rula \texttt{b et\_exit}, \texttt{run}, \texttt{ir eax}?
\begin{lstlisting}[language=x86asm]
.data
    n: .long 4
    v: .long 0x01020304, 0x05060708, 0x090a0b0c, 0x0d0e0f10
.text
.global main
main:
    mov $v, %esi
    mov $2, %ecx
    mov -8(%esi, %ecx, 4), %eax
et_exit:
    mov $1, %eax
    xor %ebx, %ebx
    int $0x80
\end{lstlisting}

\textbf{Opțiuni:}
\begin{multicols}{2}
\begin{enumerate}[label=\alph*)]
    \item 0x05060708
    \item 0x01020304 
    \item Execuție cu eroare
    \item 0x090a0b0c
\end{enumerate}
\end{multicols}
\newpage

% --- Întrebarea 7 (din 4.2) ---
\subsection*{Întrebarea 7}
Ce se va afișa pe ecran?
\begin{multicols}{2}
\begin{lstlisting}[language=x86asm]
.data
    n: .long 9
    s: .asciz "a1C95dBx3"
    t: .space 10
.text
.global main
main:
    mov $s, %esi
    mov $t, %edi
    mov $0, %ecx
    mov $0, %edx
et_loop:
    cmp n, %ecx
    je et_exit
    mov (%esi, %ecx, 1), %al
    cmp $'0', %al
    jl et2
    cmp $'9', %al
    jg et2
    mov %al, (%edi, %edx, 1)
    inc %edx
et2:
    inc %ecx
jmp et_loop
et_exit:
    movb $0, (%edi, %edx, 1)
    inc %edx
    mov $4, %eax
    mov $1, %ebx
    mov $t, %ecx
    int $0x80
    mov $1, %eax
    xor %ebx, %ebx
    int $0x80
\end{lstlisting}
\end{multicols}

\textbf{Opțiuni:}
\begin{multicols}{2}
\begin{enumerate}[label=\alph*)]
    \item aCdBx
    \item 1C95dB3
    \item 1953 
    \item a1C95dBx3
\end{enumerate}
\end{multicols}

% --- Întrebarea 8 (din 4.2) ---
\subsection*{Întrebarea 8}
Ce valori vor fi stocate în x când execuția va ajunge la eticheta \texttt{et\_exit}?
\begin{multicols}{2}
\begin{lstlisting}[language=x86asm]
.data
    x: .long 4, 2, 1, 5, 6
    n: .long 5
.text
.global main
main:
    mov $x, %esi
    mov $0, %eax
et_loop:
    cmp n, %eax
    je et_exit
    mov (%esi, %eax, 4), %ecx
    mov $1, %ebx
    sal %cl, %ebx
    mov %ebx, (%esi, %eax, 4)
    inc %eax
    jmp et_loop
et_exit:
    mov $1, %eax
    xor %ebx, %ebx
    int $0x80
\end{lstlisting}
\end{multicols}

\textbf{Opțiuni:}
\begin{multicols}{2}
\begin{enumerate}[label=\alph*)]
    \item 16, 4, 1, 25, 36
    \item 1, 2, 3, 4, 5
    \item 16, 4, 2, 32, 64 
    \item 0, 0, 0, 0, 0
\end{enumerate}
\end{multicols}
\newpage

% --- Întrebarea 9 (din 4.2) ---
\subsection*{Întrebarea 9}
Fie următorul program. Ce valori se vor regăsi în registrul \%ebx la trecerea prin eticheta \texttt{et}?
\begin{lstlisting}[language=x86asm]
.data
    v: .long 15, 21, 30, 16, 18
    n: .long 4
.text
.global main
main:
    mov $n, %esi
    mov $0, %eax
    sub n, %eax
    mov $0, %ecx
et_loop:
    cmp %eax, %ecx
    je et_exit
    mov (%esi, %ecx, 4), %ebx
et:
    dec %ecx
    jmp et_loop
et_exit:
    mov $1, %eax
    xor %ebx, %ebx
    int $0x80
\end{lstlisting}

\textbf{Opțiuni:}
\begin{multicols}{2}
\begin{enumerate}[label=\alph*)]
    \item 4, 18, 16, 30 
    \item 15, 21, 30, 16
    \item 18, 16, 30, 21
    \item 15, 21, 30, 16, 18
\end{enumerate}
\end{multicols}

% --- Întrebarea 10 (din 4.2) ---
\subsection*{Întrebarea 10}
Fie următorul program. Ce valoare va avea elementul din mijloc din vector dacă vom rula \texttt{b et\_exit}, \texttt{run}, \texttt{x/3xw 8v}?
\begin{lstlisting}[language=x86asm]
.data
    v: .long 0x01020304, 0x05060708, 0x090a0b0c
.text
.global main
main:
    mov $v, %esi
    mov $2, %ecx
    mov (%esi, %ecx, 1), %eax
    mov %eax, 4(%esi, %ecx, 1)
et_exit:
    mov $1, %eax
    xor %ebx, %ebx
    int $0x80
\end{lstlisting}

\textbf{Opțiuni:}
\begin{multicols}{2}
\begin{enumerate}[label=\alph*)]
    \item 0x0506070c
    \item 0x090a0b0c
    \item 0x01020708 
    \item 0x05060708
\end{enumerate}
\end{multicols}
\newpage

% --- Întrebarea 11 (din 4.2 / 4.3) ---
\subsection*{Întrebarea 11}
Ce valoare va fi stocată în s când execuția va ajunge la eticheta \texttt{et\_exit}?
\begin{multicols}{2}
\begin{lstlisting}[language=x86asm]
.data
    v: .long 15, 21, 30, 16, 18, 12
    n: .long 6
    s: .long 0
.text
.global main
main:
    mov $v, %esi
    mov n, %eax
    shr $1, %eax
    mov $0, %ecx

et_loop:
    cmp %eax, %ecx
    jge et_exit
    mov 0(%esi), %ebx  # (sau movl (%esi), %ebx)
    add %ebx, s        # (sau addl %ebx, s)
    add $8, %esi       # (sau addl $8, %esi)
    inc %ecx
    jmp et_loop

et_exit:
    mov $1, %eax
    xor %ebx, %ebx
    int $0x80
\end{lstlisting}
\end{multicols}

\textbf{Opțiuni:}
\begin{multicols}{2}
\begin{enumerate}[label=\alph*)]
    \item 63
    \item 66
    \item 129
    \item 23
\end{enumerate}
\end{multicols}

% --- Întrebarea 12 (din 4.3) ---
\subsection*{Întrebarea 12}
Fie următorul program. Este acesta scris corect? Daca da, de ce, daca nu, de ce?
\begin{multicols}{2}
\begin{lstlisting}[language=x86asm]
.data
    x1: .long 5
    x2: .long 12
    x3: .long 27
    n: .long 3
    s: .long 0
    formatPrintf: .asciz "%d\n"
.text
.global main
main:
    movl $x1, %edi
    movl $0, %ecx

et_loop:
    cmp n, %ecx
    je et_exit
    movl (%edi, %ecx, 4), %eax
    addl %eax, s
    incl %ecx
    jmp et_loop
et_exit:
    push s
    push $formatPrintf
    call printf
    pop %ebx
    pop %ebx

    movl $1, %eax
    movl $0, %ebx
    int $0x80
\end{lstlisting}
\end{multicols}

\textbf{Opțiuni:}
\begin{enumerate}[label=\alph*)]
    \item este scris corect, deoarece x1 e doar o adresa, ca numele unui vector.
    \item nu este scris corect, din cauza ca intram in zone de memorie;
    \item nu este scris corect, deoarece n trebuia sa fie egal cu 1;
    \item nu este scris corect, deoarece foloseste x1 pe post de array;
    \item este scris corect, dar nu va functiona conform asteptarilor;
\end{enumerate}
\newpage

% --- Întrebarea 13 (din 4.3) ---
\subsection*{Întrebarea 13}
Știind ca y este un \texttt{.long 0} declarat in sectiunea \texttt{.data}, care dintre urmatoarele poate fi valoarea initiala din x, stiind ca in urma apelului printf se va afisa la STDOUT valoarea 1572?

\begin{multicols}{2}
\begin{lstlisting}[language=x86asm]
main:
    movl $0, %ecx
    movl $10, %edi
et_while:
    cmp x, %ecx
    je et_exit
    movl x, %eax
    movl $0, %edx

    div %edi
    movl %eax, x
    movl %edx, %esi

    movl y, %eax
    mul %edi
    add %esi, %eax
    movl %eax, y
    jmp et_while

et_exit:
    push y
    push $formatStr
    call printf
    popl %ebx
    popl %ebx
\end{lstlisting}
\end{multicols}

\textbf{Opțiuni:}
\begin{multicols}{2}
\begin{enumerate}[label=\alph*)]
    \item 15720
    \item 27510 
    \item 83412
    \item 1024
    \item 1572
\end{enumerate}
\end{multicols}

\newpage
\section{Rezolvări} \label{sec:rezolvari}
\begin{multicols}{3}
\hyperref[sec:lab2]{\textbf{Test Laborator 2.1 / 2.2 / 2.3}}
\begin{enumerate}
    \item b) % 1026
    \item a) % 1
    \item a) / b) % abc12 / abc1
    \item b) / c) % 04030205 / 04030501
    \item b) % 2
    \item b) % 1234a
    \item c) % ... i r ch
\end{enumerate}
\vfill\null\columnbreak
\hyperref[sec:lab3]{\textbf{Test Laborator 3.1 / 3.2}}
\begin{enumerate}
    \item a) % 1
    \item b) % eax=0, edx=2
    \item c) % eax=3, edx=2
    \item a) % 5
    \item c) % 5
    \item a) % ety, etz, ety, ett
    \item b) % 2aaaaaad, 3
    \item a) % eax=0, edx=4
    \item b) % eax=2, edx=7
    \item b) % 4
    \item b) % etb, eth, eto, etd
    \item c) % este un ciclu infinit
    \item a) % 4
\end{enumerate}
\vfill\null\columnbreak
\hyperref[sec:lab4]{\textbf{Test Laborator 4.1 / 4.2 / 4.3}}
\begin{enumerate}
    \item b) % 11..16
    \item a) % ababaabbab
    \item b) % ABC
    \item d) % dcb
    \item a) % 05...c
    \item b) % 01...4
    \item c) % 1953
    \item c) % 16, 4, 2, 32, 64
    \item a) % 4, 18, 16, 30
    \item c) % 01020708
    \item a) % 63
    \item a) % corect (adresa)
    \item b) % 27510
\end{enumerate}

\end{multicols}

% \end{multicols}
\end{document}